Let c be the mode and \[f(x) = x^{a-1} (1-x)^{b-1},\]where \(0<x<1\).

Taking logarithms and then finding the derivative,
\begin{flalign*}
    \frac{d}{dx} log(f(x)) &= \frac{d}{dx} (a-1)log(x) + (b-1)log(1-x)\\
    &= \frac{a-1}{x} - \frac{b-1}{1-x}
\end{flalign*}
Now we can find the x-coordinate of the turning point t by finding x when \(\frac{d}{dx} log(f(x)) = 0\),
\begin{flalign*}
    \frac{a-1}{t} - \frac{b-1}{1-t} &= 0\\
    (1-t)(a-1) - (b-1)t &= 0 \\
    (1-t)(a-1) &= (b-1)t \\
    a-1 &= (b-1)t + (a-1)t \\
    a-1 &= (a+b-2)t\\
    t &= \frac{a-1}{a+b-2}  
\end{flalign*}
Now we need to determine the nature of the turning point. Finding the second derivative of \(f(x)\),
\begin{flalign*}
    \frac{d^2}{dx^2} log(f(x)) &= \frac{d}{dx} \frac{a-1}{x} - \frac{b-1}{1-x}\\
    &= -\frac{a-1}{x^2} - \frac{b-1}{(1-x)^2} \\
    &= - (\frac{a-1}{x^2} + \frac{b-1}{(1-x)^2}) 
\end{flalign*}
since \(0 < x < 1\), \(x^2 > 0\) and \((1-x)^2 > 0\) for all x,  taking \(a > 1\) and \(b > 1\) aswell,
\begin{flalign*}
    \frac{a-1}{x^2} & > 0\\
    \frac{b-1}{(1-x)^2} &> 0 
\end{flalign*}
therefore, \[- (\frac{a-1}{x^2} + \frac{b-1}{(1-x)^2}) = \frac{d^2}{dx^2} log(f(x))  < 0\]
Thus, \(t\) is the maxima and so, \[c = t = \frac{a-1}{a+b-2}\]
